\section{The Fear of That Day, and Hope}\label{sec:fear}

Among the \latin{instrumenta bonorum operum}, the tools of good works that St Benedict sets before his monks in the fourth chapter of his Rule stand three short sentences in a row: \latin{Diem iudicii timere. Gehennam expavescere. Vitam aeternam omni concupiscentia spiritali desiderare.}\endnote{\emph{Regula Benedicti} 4.44--46, Latin text as presented on Latin Wikisource (accessed 16 October 2026).} To fear the day of judgment; to dread hell; to desire eternal life with all the longing of the spirit. The three belong together, and their order is instructive. Fear comes first, and it is fear of a day and of a place; but the sentence that completes them is desire. The fear the Church teaches is a fear that turns the soul toward what it loves.

The Catechism names this fear and joins it to hope in the same breath. The message of the last judgment ``calls men to conversion while God is still giving them `the acceptable time, \ldots\ the day of salvation.' It inspires a holy fear of God and commits them to the justice of the Kingdom of God.''\endnote{\emph{Catechism of the Catholic Church} 1041, official English.} And it proclaims, the Catechism continues, the ``blessed hope'' of the Lord's return (Tit 2:13). Holy fear and blessed hope are not two moods between which the Christian must choose. They are two movements of one heart before one Judge, and the Scriptures command both: ``with fear and trembling work out your salvation'' (Phil 2:12); ``look up and lift up your heads, because your redemption is at hand'' (Luke 21:28).

\subsection{The day and the hour}

``But of that day and hour no one knoweth: no, not the angels of heaven, but the Father alone'' (Matt 24:36). The Lord who will judge has kept from his Church the date of his coming, and the Catechism repeats the reservation without softening it: ``Only the Father knows the day and the hour; only he determines the moment of its coming.''\endnote{\emph{Catechism of the Catholic Church} 1040.} Nor does anyone know the hour of his own death, which is the hour of his particular judgment. The two ignorances are the same ignorance in two forms, and Christ gives one command for both: ``Watch ye therefore, because you know not what hour your Lord will come. \ldots\ Wherefore be you also ready, because at what hour you know not the Son of man will come'' (Matt 24:42, 44).

The \emph{Imitation of Christ} turns this watchfulness into a habit of the mind:

\begin{quote}
In all that thou doest, remember the end, and how thou wilt stand before a strict judge, from whom nothing is hid, who is not bribed with gifts, nor accepteth excuses, but will judge righteous judgment. O most miserable and foolish sinner, who art sometimes in fear of the countenance of an angry man, what wilt thou answer to God, who knoweth all thy misdeeds? Why dost thou not provide for thyself against the day of judgment, when no man shall be able to be excused or defended by means of another, but each one shall bear his burden himself alone? Now doth thy labour bring forth fruit, now is thy weeping acceptable, thy groaning heard, thy sorrow well pleasing to God, and cleansing to thy soul.\endnote{\emph{The Imitation of Christ}, book~I, chapter~24, 1 (Benham).}
\end{quote}

The passage turns twice on the word \emph{now}. The fear of the judgment is not meant to paralyse but to make the present precious. ``Be thou anxious now and sorrowful for thy sins,'' the same chapter continues, ``that in the day of judgment thou mayest have boldness with the blessed.''\endnote{\emph{Imitation}, book~I, chapter~24, 4 (Benham).}

\subsection{Choosing as at the judgment}

Ignatius of Loyola made the thought of the judgment an instrument of choice. Among the rules he gives in the \emph{Spiritual Exercises} for making a good and sound election, when a soul must decide how to order its life, the fourth is this:

\begin{quote}
The fourth, looking and considering how I shall find myself on the Day of Judgment, to think how I would then want to have deliberated about the present matter, and to take now the rule which I would then wish to have kept, in order that I may then find myself in entire pleasure and joy.\endnote{Ignatius of Loyola, \emph{Spiritual Exercises}, Second Week, ``The Second Way to Make a Good and Sound Election,'' Fourth Rule, trans.\ Elder Mullan (New York: P.~J. Kenedy and Sons, 1914), pp.~92--93.}
\end{quote}

The rule asks nothing exotic. It asks the soul to borrow, for the present decision, the clarity it will have at the judgment, when every choice will be seen at last for what it was. Ignatius does not offer it to frighten but to free: the soul that chooses as it will wish to have chosen will meet the judgment ``in entire pleasure and joy.'' A Christian can make such a rule a daily practice. Before a decision, before a word spoken in anger, before a pleasure that the conscience questions, he can ask how it will look when the book is opened, and choose accordingly.

\subsection{Perfect love casteth out fear}

Fear is not the last word, and the saints never meant it to be. ``In this is the charity of God perfected with us, that we may have confidence in the day of judgment. \ldots\ Fear is not in charity: but perfect charity casteth out fear, because fear hath pain'' (1~John 4:17--18). Benedict~XVI, asking whether the image of the last judgment is not a frightening one, answers: ``I would say: it is an image that evokes responsibility, an image, therefore, of that fear of which Saint Hilary spoke when he said that all our fear has its place in love.'' And he sets the whole doctrine on its true foundation: ``The image of the Last Judgement is not primarily an image of terror, but an image of hope; for us it may even be the decisive image of hope.''\endnote{\emph{Spe salvi} 44, official English.} Justice and grace meet in the Judge. ``If it were merely grace, making all earthly things cease to matter, God would still owe us an answer to the question about justice \ldots\ If it were merely justice, in the end it could bring only fear to us all.''\endnote{\emph{Spe salvi} 47.} Because he is both just and merciful, the Christian can fear him and run to him at once.

One judgment the Christian is forbidden to make: the judgment of another's soul. ``Although we can judge that an act is in itself a grave offense,'' the Catechism teaches, ``we must entrust judgment of persons to the justice and mercy of God.''\endnote{\emph{Catechism of the Catholic Church} 1861, official English.} The only soul whose judgment a Christian is commanded to anticipate is his own.

\subsection{Verdict and vindication}

The year ends and the year begins with the same coming before the Church's eyes. On the last Sunday she hears that heaven and earth shall pass; on the first Sunday of Advent she hears that her redemption is at hand. Between those two Gospels lies the whole Christian life, lived under the eyes of a Judge who has already come once to save and will come again to judge, and whose last word in the Scriptures is a promise: ``Surely, I come quickly''. The Church's answer is the last prayer of the Bible: ``Amen. Come, Lord Jesus'' (Apoc 22:20).

A man's own case closes when he dies; what he did goes on working in the world he has left until history itself is over. God has so ordered the two judgments, and the Church confesses both: that a man is answerable at the instant of his death, without extension or appeal, and that he is more than the instant of his answering. The last day is not a rehearing. It is the moment at which the sentence already passed on a soul becomes the truth made visible about a life---the whole of it, including the part that outlived him---before the Judge who died for him, in the sight of all the saints, to the glory of God.
